%% ============================================================
%% 分片 03：原文第 2 节「${\EUN}=4$ Super Yang-Mills Theory And $S$-Duality」（原文 glshort.tex 起于第 642 行）
%% 由 tools/harvmac2tex.py 从 arXiv 源码（hep-th/0604151v3）自动生成。
%% 只译英文散文；公式、编号、\label、\cite 一律不改。
%% ============================================================

\subsection{$S$-对偶}

\label{sec:reviewsd}

对 $S$-对偶的全面回顾不幸超出了本文的范围。我们只提及几个相关的事实。

由于虚部为正，$\tau=\theta/2\pi +4\pi i/e^2$ 定义了上半平面 ${\EUH}$ 中的一个点。群
$SL(2,\Bbb{R})$ 按标准方式 $\tau\to (a\tau+b)/(c\tau+d)$ 作用在 ${\EUH}$ 上，其中 $ad-bc=1$。\footnote{忠实作用于 ${\EUH}$ 的群是商群 $PSL(2,\Bbb{R})$，但在应用于四维规范理论时，人们实际需要的是它的双重覆盖。}变换 $T:\tau\to\tau+1$ 不过是角 $\theta$ 的 $2\pi$ 平移，因此对任何规范群 $G$，它都是 $\Bbb{R}^4$ 上理论的经典对称性。$S$-对偶猜想断言，存在一个额外的量子对称性，它反演 $\tau$、把 $G$ 换成 $^L\negthinspace G$，并交换电、磁荷。此外，这个我们将称之为 $S$ 的对称性与经典对称性 $T:\tau\to\tau+1$ 结合，生成 $SL(2,\Bbb{R})$ 的一个无限离散子群 $\Gamma$。

最熟悉的情形是 $G$ 为单连通型。此时
% 式 (2.16)
\setcounter{equation}{15}
\begin{equation}\label{eq:bivolo}
S=\left(\begin{matrix} 0 & 1\\ -1 & 0\\\end{matrix}\right)
\end{equation}
作用为 $\tau\to -1/\tau$，并与 $T$ 一起生成群
$SL(2,\Bbb{Z})$。

如果 $G$ 不是单连通型，那么 $S$-变换就不是 $\tau\to -1/\tau$。相反，它是
$\tau\to -1/n_{\frak g}\tau$，其中 $n_{\frak g}$ 对 $F_4$ 为 $2$、对 $G_2$ 为 3。这个
变换可以由如下 $SL(2,\Bbb{R})$ 变换来实现：
% [WARN] 公式 \eqn\eqref{eq:bivolo} 无编号或重复
\begin{equation*}
\tag*{(2.17)}
S=\left(\begin{matrix} 0 & 1/\sqrt {n_{\frak
g}}\\ -\sqrt {n_{\frak g}} & 0\\\end{matrix}\right).
\end{equation*}
因子 2 或 3 出现在 $S$ 对 $\tau$ 的作用中，这一点
可以通过考察电、磁荷的 BPS 质量公式看出 \cite{ref46,ref47}，
并且反映了这些群的根与余根之间的关系。它也可以从非单连通型
群的 $S$-对偶的弦论解释中提取出来 \cite{ref48}。因子 2 或 3 意味着 $G_2$ 或 $F_4$ 的
对偶群不是 $SL(2,\Bbb{Z})$ 本身，而是 $SL(2,\Bbb{R})$ 的某些无限离散子群，即所谓的 Hecke 群。

其余的单李群 $Sp(k)$ 与 $Spin(2k+1)$（以及它们各自对 $\Bbb{Z}_2$ 的商）
需要一些特别的说明，因为它们是仅有的满足 $G$ 与 $^L\negthinspace G$ 的李代数不同构的单李群。
杨-米尔斯拉格朗日量、因而 $\tau$ 参数的定义，只依赖于规范群
的李代数 $\frak g$。因此，只要 $G$ 与 $^L\negthinspace G$ 具有相同的
李代数，我们就可以讨论 $S$-变换如何作用于 $\tau$，而不必区分心里想的是
规范群为 $G$ 的理论的 $\tau$，还是规范群为 $^L\negthinspace G$ 的理论的
$\tau$。这确实就是我们到目前为止隐含地所做的。

然而，对 $Sp(k)$ 与 $Spin(2k+1)$ 这一对群，没有
同样直接的方式来比较这两个 $\tau$ 参数。因此，
人们可以引入各自的规范耦合参数，比如说，对 $Sp(k)$ 或 $Sp(k)/\Bbb{Z}_2$ 规范理论取 $\tau$，
对 $Spin(2k+1)$ 或 $Spin(2k+1)/\Bbb{Z}_2=SO(2k+1)$ 规范理论取 $\tau'$。如果
人们归一化各自的 $\tau$ 参数，使得各自的 $T$-变换按 $\tau\to \tau+1$ 与
$\tau'\to\tau'+1$ 作用，那么 $S$ 按 $\tau=-1/2\tau'$、$\tau'=
-1/2\tau$ 作用，正如对 $F_4$ 那样。（因此我们对
$Spin(2k+1)$ 与 $Sp(k)$ 取 $n_{\frak g}=2$。）

这种归一化是自然的，因为它给出任意规范群下最统一的
规范理论公式。一个小小的复杂之处在于，对同样的这些群
还存在第二种也可能被认为是自然的归一化。$Sp(k)/Spin(2k+1)$ 对偶的一个
有用的弦论实现，涉及 orientifold 三平面 \cite{ref49}，实际上
提示了一种不同的归一化。在这种归一化中，
$Sp(k)$ 理论由耦合参数
$\tilde\tau=2\tau$ 描述，于是 $S$ 简单地按 $\tilde \tau =
-1/\tau'$ 作用，但 $T$ 改为按 $\tilde\tau\to \tilde \tau+2$ 作用。
（在 \cite{ref49} 的图 3 中，$Sp(k)$ 理论出现两次，
正是因为该文中 $Sp(k)$ 耦合参数 $\tilde\tau$ 的归一化使得 $T$-对称性按
$\tilde\tau\to\tilde\tau+2$ 作用。）

关于这一点，最后要说明的是，Montonen-Olive
对偶交换 $G$ 与 $^L\negthinspace G$ 这一断言并不完全是全部故事。在 $\Bbb{R}^4$ 上，这是一个足够好的描述，但
在一般的四维流形上，完整的故事要更复杂一些。四维流形上的 $G$-丛有一个特征
类 $\xi\in H^2(M,\pi_1(G))$，本文在
第~\ref{sec:centertop}~节中研究它。对 $\xi$ 的每个取值都可以定义
路径积分，而所得的配分函数 $Z_\xi$ 按对偶群
\cite{ref45,ref26} 的幺正表示变换。

这些内容对几何 Langlands 纲领有多重要？
基本的几何 Langlands 对偶就是按 $\tau\to -1/n_{\frak g}\tau$ 作用的变换 $S$。
它作用在典范参数 $\Psi$（在第~\ref{sec:canonpar}~节引入）上，方式为
$\Psi\to -1/n_{\frak g}\Psi$。基本的几何 Langlands
对偶涉及把 $\Psi=\infty$ 与 $\Psi=0$ 相比较。无论 $n_{\frak g}$ 取何值，
这两者都被 $S$ 交换。所以对基本的
几何 Langlands 对偶来说，$n_{\frak g}$ 的值并不太重要。在第~\ref{sec:genotics}~节中，我们将转向
几何 Langlands 对偶对任意 $\Psi$ 的推广。在那里，精确的对偶群是重要的，
因此 $n_{\frak g}$ 的值确实起作用。

\par\medskip{\it 超对称的变换}

关于 Montonen-Olive 对偶，有一个问题实际上
将在本文中扮演更重要的角色：它如何作用于
超对称？

我们无法通过直接检查
经典拉格朗日量来回答这个问题，因为除了由 $\tau\to \tau+1$ 生成的子群
之外，$\Gamma$ 并不由经典理论的对称性
组成。毕竟，这正是使 $S$-对偶有趣的原因。所以，要得到答案将需要
更精细的推理。

在确定答案之前，让我们问一问答案在多大程度上是
唯一的。考虑由 $S$ 与 $T$ 生成的对偶群 $\Gamma$ 的元素 $\gamma=\left(\begin{matrix}a & b
\\ c & d \\\end{matrix}\right)$。这样的元素按 $\tau\to
(a\tau+b)/(c\tau+d)$ 作用于 $\tau$，并按一个自同构作用于超对称代数。
这个自同构并非唯一确定，因为它
可以与 ${\EUN}=4$ 超对称杨-米尔斯理论（在 $\tau$ 取固定值时）的一个对称性复合。一个重要的
简化是，根据 Montonen-Olive 猜想，$\gamma$ 与 Poincar\'e 群交换。此外，
人们可以定义它使得它与整体 $R$-对称群
$SU(4)_{{\cal R}}$ 交换。\footnote{为说明这一点，我们首先注意到，用 $\gamma$ 作共轭
生成 $SU(4)_{{\cal R}}$ 的一个自同构。这个自同构必然是内自同构，因为
经典理论没有任何对称性既在时空上平凡地作用、又按
$SU(4)_{{\cal R}}$ 的外自同构作用。最后，既然
$\gamma$ 生成 $SU(4)_{{\cal R}}$ 的内自同构，我们
就可以把它与一个生成逆自同构的 $SU(4)_{{\cal R}}$ 元素复合，从而得到
一个与 $SU(4)_{{\cal R}}$ 交换的对偶对称性。}综合这些事实可以推出，
$\gamma$ 在按 $\bar{\bf 4}$（即 $\bar Q_{AX}$）变换的
超对称上表现为标量乘法 $\exp(i\hat\phi)$，而在按
${\bf 4}$（即 $ Q_{\dot A}^Y$）变换的超对称上表现为 $\exp(-i\hat\phi)$。此外，$\hat\phi$ 必须是
实数，以保持代数的实结构（相对于这个实结构，在 Lorentz 号差下 $P$ 是厄米的，
而 $\bar Q$ 是 $Q$ 的厄米伴随）。我们将称这些对称性为
$U(1)$ 手征转动。$\gamma$ 的作用定义到
$SU(4)_{{\cal R}}$ 的中心的一个元素为止。该中心由
元素 ${\EUI}$ 生成，它在 $SU(4)_{{\cal R}}$ 的 ${\bf 4}$ 上作用为 $i$
、在 $\bf{ \bar 4}$ 上作用为 $-i$。因此
$\hat\phi(\gamma)$ 定义到
$\hat\phi(\gamma)\to\hat\phi(\gamma)+\pi/2$ 为止。

为确定 $\hat\phi(\gamma)$，我们可以在任意方便的态中
计算 $\gamma$ 的作用。我们选择在理论的 Coulomb 分支上进行计算，在那里规范
群破缺为 $U(1)^r$，而超对称代数被电、磁荷中心扩张。对
$U(1)$ 或 $U(1)^r$ 这样的阿贝尔规范群，人们可以显式地
算出一切，从而确定 $\gamma$ 如何作用于超对称。或者，人们可以利用
${\cal N}=4$ 超对称杨-米尔斯理论作为 IIB 型超弦理论中 $D3$-膜上
规范理论的实现；这给出一个几何的方式来
对 $G=U(1)$ 确定 $\gamma$ 对超荷的作用（而且
这一方法也直接推广到 $U(N)$ 这样的经典群）。


我们将采用另一种方法来确定结果，
即考察 BPS 态的质量公式。为简化
记号，我们聚焦于超对称代数的一个 ${\EUN}=2$ 子代数，它有一对右手
超荷 $Q_{\dot A}^i$，$i=1,2$，以及单个中心荷
$Z$。它们满足 % 式 (2.18)
\setcounter{equation}{17}
\begin{equation}\label{eq:hilj}
\{Q_{\dot A}^i,Q_{\dot
B}^j\}=\eps_{\dot A\dot B}\eps^{ij} Z.
\end{equation}
作为无质量规范场的 ${\EUN}=2$ 超伴侣的复标量
取值于李代数 ${\frak g}$ 的 Cartan 子代数 ${\frak t}$。我们归一化这个标量，使其动能项
与 $\tau$ 无关，并记它的期望值为
$\vec{\phi}$。电、磁荷记为
$\vec{\eurm n}$ 与 $\vec{\eurm m},$ 取值于 $\frak g$ 的权格与
余权格。向量 $\vec{\phi},\vec{\eurm
n},\vec{\eurm m}$ 定义到 Weyl 群的作用为止。

为使我们的约定清晰，我们将给出
$G=U(1)$ 时 $\eurm n$ 与 $\eurm m$ 的显式定义。物理上，
人们通常用曲率为 $F'=dA'$ 的实联络 $A'$ 来描述 $U(1)$ 规范理论。
我们把 Minkowski 时空写成
$\Bbb{R}^{1,3}=\Bbb{R}\times \Bbb{R}^3$，其中第一个因子
参数化时间，第二个因子参数化空间。为了与 Maxwell 方程一致（例如见
\cite{ref51} 第 42 页），电、磁场定义为
$E_i=F'{}^{0i}$ 与 $B_i=\half\epsilon_{ijk}F'_{jk}$（这里
$\epsilon_{ijk}$ 是 $\Bbb{R}^3$ 上的反对称张量，且
$\epsilon_{123}=1$）。于是 $\eurm n$ 与 $\eurm m$ 定义为
% 式 (2.19)
\setcounter{equation}{18}
\begin{equation}\label{eq:eltro}
\begin{aligned}{\eurm n} & = {1\over 4\pi}\int_Y d\Omega\,
\vec n\cdot \vec E\\ {\eurm m}&={1\over 4\pi}\int_Yd\Omega\,\vec
n\cdot \vec B\\\end{aligned}
\end{equation}
其中 $Y$ 是 $\Bbb{R}^3$ 中无穷远处的一个大球面，
其体积形式为 $d\Omega$，法向量为 $\vec n$。
改用与 $A'$ 通过 $A=-iA'$ 与 $F=-iF'$ 相联系的反厄米规范场 $A$ 来表示
（回忆式 \eqref{eq:boloko}\ 与 \eqref{eq:obolko}），我们有
% 式 (2.20)
\setcounter{equation}{19}
\begin{equation}\label{eq:feltro}
\begin{aligned}{\eurm n} & = {i\over 4\pi}\int_Y d\Omega\,
n^i F^{0i}\\ {\eurm m}&={i\over 4\pi}\int_Yd\Omega\,\half
\epsilon_{ijk}n^i F^{jk}.\\\end{aligned}
\end{equation}
用微分形式来表示，第二个
公式为 % 式 (2.21)
\setcounter{equation}{20}
\begin{equation}\label{eq:ofeltro}
{\eurm m}={i\over 2\pi}\int_YF
=\int_Y c_1({\cal L}),
\end{equation}
其中 ${\cal L}$ 是 $A$ 为其上一个联络的
线丛。

中心荷 $Z$ 由下式给出 % 式 (2.22)
\setcounter{equation}{21}
\begin{equation}\label{eq:pikf}
Z=\sqrt{2\over {\rm
Im}\,\tau}\,\left(\vec{\eurm m}~\vec{\eurm n}\right)\cdot
\left(\begin{matrix}\tau\vec \phi \\\vec \phi \\\end{matrix}\right)=\sqrt {2\over
{\rm Im}\,\tau}\,\vec \phi\cdot (\vec{\eurm n}+\tau \vec{\eurm m})
\end{equation}
对于单连通型 $G$，Coulomb 分支上的低能阿贝尔规范理论
允许有一个 $SL(2,\Bbb{Z})$ 对偶对称群，使得一般元素
$\eurm M=\left(\begin{matrix}a&b\\ c&d\\\end{matrix}\right)$ 按
$\tau\to (a\tau+b)/(c\tau+d)$ 作用于 $\tau$，并按如下方式作用于 $\vec{\phi},\,
\vec{{\eurm m}},\, \vec{{\eurm n}} $：
% 式 (2.23)
\setcounter{equation}{22}
\begin{equation}\label{eq:xilj}
\begin{aligned}\vec{\phi} & \to
\vec{\phi}\\\left(\begin{matrix}\vec{{\eurm m}} & \vec{{\eurm
n}}\end{matrix}\right)&\to \left(\begin{matrix}\vec{{\eurm m}} & \vec{{\eurm
n}}\end{matrix}\right){\eurm M}^{-1}=\left(\begin{matrix} \vec{{\eurm m}} &
\vec{{\eurm n}}\end{matrix}\right)\left(\begin{matrix}d & -b  \\ -c & a
\\\end{matrix}\right).\end{aligned}
\end{equation}
$S$-对偶猜想断言，低能理论的这个
对称性实际上延拓为整个理论的对称性，并且必然以同样的方式作用。对偶群 $\Gamma$ 的生成元
$T:\tau\to\tau+1$ 延拓到整个理论是清楚的（因为它只是 $\theta$ 角的 $2\pi$ 平移），
所以猜想真正说的是 $S:\tau\to
-1/\tau$ 也同样延拓到整个理论。$T$ 在 \eqref{eq:xilj}\ 中所断言的作用
可以在整个理论中通过直接计算看出来 \cite{ref50}。这个直接计算对
所有 $G$ 都成立，无论是否单连通型；它的一个变体将在
第~\ref{sec:tophooft}~节中给出。由 \eqref{eq:pikf}\ 与 \eqref{eq:xilj}，可以推出
中心荷的变换： % 式 (2.24)
\setcounter{equation}{23}
\begin{equation}\label{eq:mujk}
Z\to
{|c\tau+d|\over {c\tau+d}}\, Z.
\end{equation}
那么，中心扩张的
超对称代数意味着右手
超对称乘以
% 式 (2.25)
\setcounter{equation}{24}
\begin{equation}\label{eq:tokio}
\exp(-i\hat\phi(\gamma))=\left({|c\tau+d|\over
{c\tau+d}}\right)^{1/2}.
\end{equation}

这个平方根意味着，作用在
超对称上的群是对偶群 $\Gamma$ 的双重覆盖。
这个扩张由一个二阶对称性 $(-1)^F$ 给出，它改变
所有费米子的符号。这个对称性是 $SU(4)_{\cal R}$ 的中心生成元 ${\EUI} $ 的{\it 平方}。我们之所以避免
让 $\Gamma$ 被整个中心扩张，是因为我们把
$\vec{\phi}$ 取为在 $\Gamma$ 下不变，而不是只不变到一个符号。

当 $G$ 不是单连通型时，这个论证需要更仔细地
表述。对偶群由 $T:\tau\to
\tau+1$ 与 $S:\tau\to -1/n_{\frak g}\tau$ 生成。$T$ 使
超对称保持不变，所以对 $T$ 有 $\hat\phi=0$，这与 \eqref{eq:tokio}
相符。至于 $S$，其作用现在为 % 式 (2.26)
\setcounter{equation}{25}
\begin{equation}\label{eq:xiljo}
\begin{aligned}\vec{\phi} &
\mapsto {\frak R}\cdot\vec{\phi}\\\left(\begin{matrix}\vec{{\eurm m}} &
\vec{{\eurm n}}\end{matrix}\right)&\mapsto \left(\begin{matrix}{\frak R}\cdot
\vec{{\eurm m}} & {\frak R}\cdot\vec{{\eurm
n}}\end{matrix}\right)\left(\begin{matrix}0 & -1/\sqrt n_{\frak g} \\ \sqrt
n_{\frak g} & 0 \\\end{matrix}\right).\end{aligned}
\end{equation}
这里 $\frak R$ 是 $\frak t$ 的一个正交
变换，对单连通型群它可以取为
平凡的，对非单连通型群则在 \cite{ref47} 中描述。$\frak R$ 的平方属于 Weyl 群，所以 $S$
的平方在模空间上平凡地作用。$\frak R$ 不
影响中心荷变换的计算，而一个小计算表明 $n_{\frak g}$ 会抵消，
所以 \eqref{eq:tokio}\ 在一般情形下仍然成立。
\begin{yizhu}[易错点]
按：「单连通型」（simply-laced）这一译名容易与拓扑上的「单连通」（$\pi_1(G)=0$）混淆，两者无关。simply-laced 指 Dynkin 图上没有多重边、所有根等长，即 $A_n$、$D_n$、$E_6$、$E_7$、$E_8$；与之相对的是 $B_n$、$C_n$、$F_4$、$G_2$（有长根与短根）。本文用 $n_{\frak g}$ 统一记这一「链环数」：单连通型为 $1$，$B_n$、$C_n$、$F_4$ 为 $2$，$G_2$ 为 $3$。译名以术语表为准。
\end{yizhu}

\begin{yizhu}[含糊点拨]
按：正文这里只说 $n_{\frak g}$ 对 $F_4$ 为 $2$、对 $G_2$ 为 $3$，似乎只涉及例外群；但 $B_n$、$C_n$（即 $Spin(2n+1)$ 与 $Sp(n)$）同样有长根与短根，$n_{\frak g}$ 也等于 $2$。正文紧接着就用「因此我们对 $Spin(2k+1)$ 与 $Sp(k)$ 取 $n_{\frak g}=2$」把这一点补上了，读时不要以为 $n_{\frak g}=2$ 只适用于 $F_4$。这个数就是 Frenkel 综述（\emph{Lectures on the Langlands Program and Conformal Field Theory}）第 4 节所说的 lacing number。
\end{yizhu}

\begin{yizhu}[含糊点拨]
按：正文写「$S$ 按 $\tau=-1/2\tau'$ 作用」，字面上用了等号，这里应读作映射 $\tau\mapsto -1/(2\tau')$（以及对 $\tau'$ 的相应映射），即先把 $Sp(k)$ 与 $Spin(2k+1)$ 的耦合互换，再做 $n_{\frak g}=2$ 的 $S$-变换。验证：矩阵 $\left(\begin{matrix}0 & 1/\sqrt{n}\\ -\sqrt n & 0\\\end{matrix}\right)$ 对 $\tau$ 的作用是 $(a\tau+b)/(c\tau+d)=-\dfrac{1/\sqrt n}{\sqrt n\,\tau}=-\dfrac{1}{n\tau}$，取 $n=2$ 并交换 $\tau$ 与 $\tau'$ 即得。原文把「作用」写成等号的这种写法在下文还会出现，宜按映射理解。
\end{yizhu}

\begin{yizhu}[物理图像]
按：式~\eqref{eq:tokio} 中的平方根不是书写上的方便，而有物理含义：中心荷按 $|c\tau+d|/(c\tau+d)$ 变换，这个相位是模 $2\pi$ 的；而超荷是中心荷的「平方根」（BPS 质量 $M=|Z|$），所以超荷拿到的相位是它的一半。取遍所有 $\gamma$，这一半相位一般不能拼成 $\Gamma$ 的表示，只能拼成 $\Gamma$ 的双重覆盖的表示；多出来的 $\Bbb{Z}_2$ 正是 $(-1)^F$，它把超荷变号而对玻色子无作用。正文随后说明 $(-1)^F$ 是 $SU(4)_{\cal R}$ 的中心元 ${\EUI}$ 的平方，因此这个扩张已经包含在 $R$-对称群的中心之内，而不是新东西。
\end{yizhu}

\begin{yizhu}[含糊点拨]
按：为什么 $\hat\phi(\gamma)$ 只能定义到 $\pi/2$ 而不能再细？因为 $\gamma$ 与 $SU(4)_{\cal R}$ 的中心元素复合后仍是同一个对偶对称性。中心元 ${\EUI}$ 在 ${\bf 4}$ 上作用为 $i$、在 $\bar{\bf 4}$ 上作用为 $-i$；把它并进 $\gamma$，相当于在 $\bar Q$ 上多乘 $-i$、在 $Q$ 上多乘 $i$，也就是把 $\hat\phi$ 平移 $\pi/2$ 的奇数倍。因此 $\hat\phi(\gamma)$ 是一个 $\Bbb{R}/\frac{\pi}{2}\Bbb{Z}$ 中的元素，只有当群的中心平凡、或进一步固定规范时，才能把它提升为逐点确定的数。
\end{yizhu}

\begin{yizhu}[原文疑误]
按：原文此句在动词 answer 前多了一个冠词 the（「We cannot the answer this question」），所据 arXiv 源码（hep-th/0604151v3）亦然；英文本意是「We cannot answer this question」。译文按后者处理。
\end{yizhu}

\begin{yizhu}[物理图像]
按：这一段把 $S$-对偶的群论部分讲完了：$\tau$ 住在上半平面 ${\EUH}$ 中，$T:\tau\to\tau+1$ 与 $S$ 一起生成 $\Gamma\subset SL(2,\Bbb{R})$；单连通型时 $\Gamma=SL(2,\Bbb{Z})$，非单连通型时是 Hecke 群（$SL(2,\Bbb{R})$ 的指数有限的离散子群）。几何 Langlands 关心的不是 $\tau$ 的数值，而是 $\Psi$ 的两个端点：$\Psi=\infty$（电侧，Wilson 算子所在）与 $\Psi=0$（磁侧，'t Hooft 算子所在），它们恰好被 $S$ 互换；而 $\Psi$ 的一般值需要整个 $\Gamma$ 才能描述——这就是本节末句「$n_{\frak g}$ 的值确实起作用」的含义。
\end{yizhu}
